% GENERATED FILE. Edit canonical lessons, then run npm run generate:books.
% canonical-source-hash: fnv1a64:5f9ae3c8f43347b3
% canonical-lessons: FR-C185-muscle, FR-C185-pouce, FR-C185-blessure, FR-C185-cicatrice, FR-C185-rhume

\chapter{Muscle, Thumb, Wound, Scar, Cold}
\label{ch:fr-muscle-thumb-wound-scar-cold}
\hlchaptermodality{185}

\begin{chapteropening}
\textbf{By the end of this chapter:} \emph{I can say a muscle, a thumb, a wound, a scar and a cold.}

Complete the last lesson of chapter 185: I can say a muscle, a thumb, a wound, a scar and a cold.
\end{chapteropening}

\section[le muscle]{le muscle — a muscle}

\label{lesson:FR-C185-muscle}



\begin{quote}
Before the new one: say the French for to disappear, then the French for to borrow.
\end{quote}

\subsection*{You'll want to know: le muscle}
\textbf{le muscle} — \textquotedblleft{}a muscle\textquotedblright{}.

\begin{grammarlens}[title={What you've built}]
One more word for this chapter.
\end{grammarlens}

\subsection*{Your turn}
\begin{itemize}
\raggedright
  \item \emph{Say it:} \emph{le muscle}
  \item \emph{Say it:} \emph{le muscle}, once more
  \item \emph{Say it:} say \emph{le muscle}
  \item \emph{Recall:} say the French for to disappear, then the French for to borrow, then say \emph{le muscle} again
\end{itemize}

\begin{tcolorbox}[breakable,colback=teal!4,colframe=teal!35!black,title={Before you move on}]
What does le muscle mean? (\textquotedblleft{}A muscle\textquotedblright{}.) Say it once more.
\end{tcolorbox}
\section[le pouce]{le pouce — a thumb}

\label{lesson:FR-C185-pouce}



\begin{quote}
Before the new one: say the French for to borrow, then the French for a muscle.
\end{quote}

\subsection*{You'll want to know: le pouce}
\textbf{le pouce} — \textquotedblleft{}a thumb\textquotedblright{}.

\begin{grammarlens}[title={What you've built}]
One more word for this chapter.
\end{grammarlens}

\subsection*{Your turn}
\begin{itemize}
\raggedright
  \item \emph{Say it:} \emph{le pouce}
  \item \emph{Say it:} \emph{le pouce}, once more
  \item \emph{Say it:} say \emph{le pouce}
  \item \emph{Recall:} say the French for to borrow, then the French for a muscle, then say \emph{le pouce} again
\end{itemize}

\begin{tcolorbox}[breakable,colback=teal!4,colframe=teal!35!black,title={Before you move on}]
What does le pouce mean? (\textquotedblleft{}A thumb\textquotedblright{}.) Say it once more.
\end{tcolorbox}
\section[la blessure]{la blessure — a wound}

\label{lesson:FR-C185-blessure}



\begin{quote}
Before the new one: say the French for a muscle, then the French for a thumb.
\end{quote}

\subsection*{You'll want to know: la blessure}
\textbf{la blessure} — \textquotedblleft{}a wound\textquotedblright{}.

\begin{grammarlens}[title={What you've built}]
One more word for this chapter.
\end{grammarlens}

\subsection*{Your turn}
\begin{itemize}
\raggedright
  \item \emph{Say it:} \emph{la blessure}
  \item \emph{Say it:} \emph{la blessure}, once more
  \item \emph{Say it:} say \emph{la blessure}
  \item \emph{Recall:} say the French for a muscle, then the French for a thumb, then say \emph{la blessure} again
\end{itemize}

\begin{tcolorbox}[breakable,colback=teal!4,colframe=teal!35!black,title={Before you move on}]
What does la blessure mean? (\textquotedblleft{}A wound\textquotedblright{}.) Say it once more.
\end{tcolorbox}
\section[la cicatrice]{la cicatrice — a scar}

\label{lesson:FR-C185-cicatrice}



\begin{quote}
Before the new one: say the French for a thumb, then the French for a wound.
\end{quote}

\subsection*{You'll want to know: la cicatrice}
\textbf{la cicatrice} — \textquotedblleft{}a scar\textquotedblright{}.

\begin{grammarlens}[title={What you've built}]
One more word for this chapter.
\end{grammarlens}

\subsection*{Your turn}
\begin{itemize}
\raggedright
  \item \emph{Say it:} \emph{la cicatrice}
  \item \emph{Say it:} \emph{la cicatrice}, once more
  \item \emph{Say it:} say \emph{la cicatrice}
  \item \emph{Recall:} say the French for a thumb, then the French for a wound, then say \emph{la cicatrice} again
\end{itemize}

\begin{tcolorbox}[breakable,colback=teal!4,colframe=teal!35!black,title={Before you move on}]
What does la cicatrice mean? (\textquotedblleft{}A scar\textquotedblright{}.) Say it once more.
\end{tcolorbox}
\section[le rhume]{le rhume — a cold}

\label{lesson:FR-C185-rhume}



\begin{quote}
Before the new one: say the French for a wound, then the French for a scar.
\end{quote}

\subsection*{You'll want to know: le rhume}
\textbf{le rhume} — \textquotedblleft{}a cold\textquotedblright{}.

\begin{grammarlens}[title={What you've built}]
One more word for this chapter.
\end{grammarlens}

\subsection*{Your turn}
\begin{itemize}
\raggedright
  \item \emph{Say it:} \emph{le rhume}
  \item \emph{Say it:} \emph{le rhume}, once more
  \item \emph{Say it:} say \emph{le rhume}
  \item \emph{Recall:} say the French for a wound, then the French for a scar, then say \emph{le rhume} again
\end{itemize}

\begin{tcolorbox}[breakable,colback=teal!4,colframe=teal!35!black,title={Before you move on}]
What does le rhume mean? (\textquotedblleft{}A cold\textquotedblright{}.) Say it once more.
\end{tcolorbox}
